
\documentclass[12pt,a4paper]{article}

\usepackage{iftex}
\ifPDFTeX
    \usepackage[T1]{fontenc}
    \usepackage[utf8]{inputenc}
    \usepackage{lmodern}
\else
    \usepackage{fontspec}
    \setmainfont{Latin Modern Roman}
    \setsansfont{Latin Modern Sans}
    \setmonofont{Latin Modern Mono}
\fi
\usepackage[polish]{babel}
\usepackage{geometry}
\usepackage[hidelinks]{hyperref}
\usepackage{url}
\emergencystretch=2em
\widowpenalty=10000
\clubpenalty=10000

\geometry{margin=2.5cm}

\begin{document}

\begin{center}
    {\LARGE \textbf{Projekt przejściowy}}\\[1cm]
    {\large \textbf{Temat projektu}}
\end{center}

\bigskip

Wykorzystanie metod uczenia maszynowego do wykrywania anomalii
w danych pochodzących z sieci czujników IoT oraz opracowanie
systemu monitorowania, wizualizacji i raportowania wykrytych
nieprawidłowości.

\vspace{1cm}

\section*{Skład grupy}

\begin{enumerate}
    \item Mateusz Buchczyk 272816
    \item Remigiusz Szmyt 272819
\end{enumerate}

\clearpage
\section{Przegląd literatury i rozwiązań technicznych}
\label{sec:literatura}

Przegląd obejmuje cztery książki podręcznikowe, cztery publikacje naukowe
i pięć źródeł technicznych dotyczących architektury ESP32--MQTT--serwer--HMI
oraz detekcji anomalii.

\subsection{Podstawy teoretyczne i podręczniki}

Wybrane podręczniki porządkują podstawy uczenia maszynowego oraz metodykę
pracy z danymi. Stanowią punkt wyjścia do omówienia algorytmów detekcji
i uzasadnienia sposobu ich porównania w projekcie.

\subsubsection{Russell i Norvig: Artificial Intelligence: A Modern Approach}
Podręcznik Russella i Norviga przedstawia sztuczną inteligencję jako dziedzinę
obejmującą m.in. rozumowanie, podejmowanie decyzji, reprezentację wiedzy
i uczenie. Dla projektu istotne są części dotyczące uczenia na przykładach,
modeli probabilistycznych oraz niepewności~\cite{russell2020}.

Źródło wybrano jako podstawę pojęciową projektu. Pomaga wyjaśnić, dlaczego
system uczący się na pomiarach różni się od zestawu ręcznie ustalonych progów
temperatury lub oświetlenia. Jego rola jest ogólna: szczegółowy opis Isolation
Forest i One-Class SVM oparto na oryginalnych publikacjach, a nie na samym
podręczniku. Taki podział pozwala zachować właściwy zakres cytowań i nie
przypisywać książce szczegółów konkretnej implementacji.

W projekcie wiedza ta pomaga oddzielić trzy poziomy: odczyt z ESP32, ocenę
nietypowości wektora pomiarowego oraz reakcję interfejsu HMI. Wniosek projektowy
jest następujący: alarm modelu należy prezentować jako sygnał wymagający
interpretacji, a nie automatyczne rozpoznanie uszkodzenia czujnika. Zmiana
warunków otoczenia może być rzeczywista i poprawnie zmierzona, mimo że model
uzna ją za nietypową. Podręcznik daje zatem ramy do opisania celu systemu,
zakresu jego autonomii i ograniczeń wnioskowania na podstawie danych.

\subsubsection{Zhou: Machine Learning}
Książka Zhi-Hua Zhou obejmuje podstawy uczenia maszynowego, wybór i ocenę
modeli, drzewa decyzyjne, maszyny wektorów nośnych, uczenie zespołowe,
grupowanie oraz dobór cech. Oficjalny spis treści wskazuje osobne rozdziały
poświęcone SVM, ocenie modeli i selekcji cech. Są to obszary bezpośrednio
związane z planowanym porównaniem detektorów~\cite{zhou2021}. 

Pozycję wybrano, ponieważ zapewnia wspólny język opisu dwóch odmiennych
podejść: zespołu losowych drzew oraz metody wykorzystującej jądro.
Pozwala omówić reprezentację próbek, generalizację i dobór parametrów przed
przejściem do specyfiki detekcji anomalii. Rozdział o klasycznym SVM nie
zastępuje jednak wyprowadzenia wariantu jednoklasowego, ponieważ ten wymaga odrębnego
odwołania do pracy Schölkopfa i współautorów.

Konkretnym wkładem w projekt jest uporządkowanie planu porównania modeli
na tych samych sześciu cechach pomiarowych. Temperatura, wilgotność, ciśnienie,
oświetlenie i dwie miary przyspieszenia mają odmienne skale oraz znaczenie
fizyczne. Wiedza o cechach i modelach pomaga uzasadnić analizę skalowania
przed One-Class SVM oraz ocenę przydatności każdego kanału. Dobór parametrów
powinien być oddzielony od końcowego testu. Są to decyzje metodyczne projektu,
nie gotowe zalecenia książki dotyczące konkretnie ESP32 ani potwierdzenie
optymalności przyjętego podziału 70/15/15.

\subsubsection{Goodfellow, Bengio i Courville: Deep Learning}
Podręcznik omawia matematyczne podstawy uczenia maszynowego, sieci głębokie,
regularyzację, optymalizację i uczenie reprezentacji. Dla analizowanego
systemu szczególnie użyteczne są rozdziały o modelowaniu sekwencji,
metodyce praktycznej i autoenkoderach. Autoenkoder przekształca dane do
reprezentacji wewnętrznej i odtwarza wejście, a ograniczenia architektury
lub procesu uczenia mają prowadzić do uchwycenia użytecznych właściwości
danych~\cite{goodfellow2016}.

Książkę wybrano jako teoretyczne tło alternatywy dla Isolation Forest
i One-Class SVM. Umożliwia wyjaśnienie, czym różni się ręcznie przygotowany
wektor pomiarów od reprezentacji uczonej przez sieć. Nie stanowi jednak
samodzielnego dowodu, że duży błąd rekonstrukcji zawsze oznacza awarię,
ani że metoda głęboka będzie skuteczniejsza od prostszych detektorów.
Zastosowanie rekonstrukcji do anomalii czasowych rozwija osobno omówiona
publikacja Malhotry i współautorów.

Wkład w projekt polega przede wszystkim na świadomym określeniu jego zakresu.
Przy około 39 godzinach zarejestrowanych pomiarów liczba wektorów nie jest
równoznaczna z dużą różnorodnością stanów pracy. Wdrożenie autoenkodera
wymagałoby dodatkowego uzasadnienia architektury, regularyzacji i procedury
uczenia. Dlatego książka wspiera opis możliwego rozszerzenia systemu,
bez deklarowania wykonanych eksperymentów z sieciami. Dla HMI wynika stąd
potrzeba prezentowania rodzaju modelu i znaczenia jego wyniku: błąd
rekonstrukcji, wynik izolacji i funkcja decyzyjna SVM nie są tą samą wielkością
i nie powinny być bez uzasadnienia porównywane na wspólnej skali.

\subsubsection{Géron: Hands-On Machine Learning}
Trzecie wydanie książki Gérona przedstawia praktyczny przebieg projektu
uczenia maszynowego z użyciem scikit-learn, Keras i TensorFlow. Towarzyszące
mu materiały autora obejmują przygotowanie danych, klasyfikację, SVM,
drzewa, metody zespołowe i uczenie bez nadzoru. Zaletą publikacji jest
powiązanie pojęć z wykonywalnymi przykładami w notatnikach
Python~\cite{geron2022}.

Pozycję wybrano ze względu na przełożenie metodyki ML na organizację kodu
i eksperymentów. W projekcie potrzebny jest powtarzalny ciąg przekształceń
od tabeli pomiarów do predykcji, a nie jednorazowe uruchomienie algorytmu.
Przykłady nie dotyczą dokładnie badanego zestawu czujników, dlatego powinny
być traktowane jako wzorce pracy, a nie gotowe rozwiązanie problemu.
W szczególności sposób podziału danych trzeba dostosować do ich
chronologicznego charakteru.

Konkretnym zastosowaniem jest przygotowanie wspólnego procesu skalowania
i inferencji dla One-Class SVM: parametry przekształceń należy wyznaczać
na treningu, a potem stosować bez ponownego dopasowania do walidacji, testu
i nowych wiadomości z serwera MQTT. Warto utrwalić model razem z kolejnością
cech i konfiguracją przetwarzania, aby HMI pokazywał wyniki zgodne
z eksperymentem. Książka pomaga również uzasadnić rozdzielenie strojenia
od oceny końcowej. Jej wykorzystanie nie oznacza przyjęcia wszystkich
bibliotek z tytułu: TensorFlow i Keras nie są konieczne do porównania
Isolation Forest z One-Class SVM.

\subsection{Publikacje naukowe}

Poniższe publikacje przedstawiają podstawy badanych detektorów, przykład
ich zastosowania w IoT oraz metody analizy szeregów czasowych. Ich omówienie
łączy opis metod z oceną przydatności dla projektowanego systemu.

\subsubsection{Liu, Ting i Zhou: Isolation Forest}
Autorzy proponują detekcję opartą na izolowaniu obserwacji przez losowe
podziały przestrzeni cech. W drzewie wybierana jest cecha i punkt podziału,
a obserwacje rzadkie i odmienne od większości zwykle można odizolować
mniejszą liczbą podziałów. Wynik lasu wynika ze średniej długości ścieżki,
z odpowiednią normalizacją. Publikacja analizuje również próbkowanie danych
oraz koszt obliczeń i porównuje metodę z innymi detektorami, wykorzystując
m.in. AUC i czas przetwarzania~\cite{liu2008}.

Artykuł wybrano jako pierwotne źródło jednego z dwóch głównych algorytmów
projektu. Uzasadnia zainteresowanie metodą, która nie wymaga etykiety każdej
próbki ani obliczania wszystkich wzajemnych odległości. Podpróbkowanie
i zespół drzew są istotne dla analizy regularnie napływających pomiarów.
Przedstawione wyniki dotyczą jednak zbiorów autorów i nie dowodzą przewagi
na danych z trzech węzłów ESP32.

W projekcie metoda będzie oceniać wspólny wektor sześciu kanałów, co pozwala
badać nietypowe kombinacje wartości, a nie wyłącznie przekroczenia progów
pojedynczych czujników. Należy zapisywać liczbę drzew, wielkość podpróbki
i ziarno losowania, aby umożliwić odtworzenie porównania z One-Class SVM.
Przy przekazywaniu wyniku do HMI trzeba uwzględnić konwencję implementacji:
oryginalny wynik anomalii i wynik scikit-learn nie mają tej samej orientacji.
Sam detektor nie identyfikuje przyczyny zdarzenia i bez cech czasowych nie
modeluje kolejności próbek; te ograniczenia powinny towarzyszyć interpretacji
alarmów.

\subsubsection{Arnau Muñoz i współautorzy: jakość danych w infrastrukturze IoT}
Artykuł \emph{Anomaly detection system for data quality assurance in IoT
infrastructures based on machine learning} opisuje system wykrywający
podejrzane dane w heterogenicznej infrastrukturze Uniwersytetu w Alicante.
Autorzy łączą zbieranie wiadomości przez MQTT, analizę i przygotowanie cech
z detekcją wykorzystującą Isolation Forest. Celem jest ochrona jakości
danych przekazywanych aplikacjom; wykrycie nietypowego pakietu nie jest
utożsamiane ze szczegółową diagnozą awarii~\cite{arnau2024}.

Publikację wybrano jako najbliższy architektonicznie przykład zastosowania.
Pokazuje ona, że użyteczny detektor wymaga również warstwy akwizycji,
przechowywania i przygotowania danych. Autorzy analizują dane pochodzące
z wielu systemów kampusowych, dlatego ich liczby cech, skala infrastruktury
i wyniki nie mogą być bezpośrednio przenoszone na trzy węzły ESP32.

Konkretnym wkładem jest wzorzec oddzielenia odbioru telemetrii od analizy ML.
W naszym systemie jeden wiersz CSV opisuje pojedynczy kanał, więc przed
detekcją konieczne jest utworzenie spójnego wektora pomiarowego. Informacje
o urządzeniu, kanale i czasie odbioru należy zachować do późniejszej
interpretacji alarmu. W odróżnieniu od celu filtrowania strumienia przedstawionego
przez autorów proponuje się pozostawienie surowych danych i oznaczanie
podejrzanych obserwacji w HMI. Umożliwi to sprawdzenie fałszywych alarmów
i odróżnienie błędu danych od rzeczywistej zmiany warunków fizycznych.

\subsubsection{Schölkopf i współautorzy: One-Class SVM}
Praca \emph{Estimating the Support of a High-Dimensional Distribution}
dotyczy estymacji obszaru przestrzeni, w którym skupia się rozkład danych.
Autorzy konstruują funkcję o różnych znakach wewnątrz i poza wyznaczonym
obszarem, wykorzystując jądra oraz optymalizację kwadratową. Jest to
rozszerzenie idei maszyn wektorów nośnych na dane bez etykiet klas.
Publikacja zawiera opis algorytmu, analizę teoretyczną i doświadczenia
na danych sztucznych oraz rzeczywistych~\cite{scholkopf2001}.

Źródło wybrano, ponieważ wyjaśnia podstawy drugiego badanego detektora,
a nie jedynie sposób wywołania biblioteki. One-Class SVM opisuje obszar
zachowania odniesienia, co stanowi wartościowy kontrast wobec mechanizmu
izolacji obserwacji. Parametr $\nu$ wiąże się z górnym ograniczeniem udziału
błędów treningowych i dolnym ograniczeniem udziału wektorów nośnych,
przy założeniach podanych przez autorów. Nie jest gwarantowanym odsetkiem
rzeczywistych awarii ani fałszywych alarmów w przyszłych danych.

Dla projektu wkład obejmuje uzasadnienie uczenia na możliwie wiarygodnym
fragmencie normalnej pracy i kontroli zanieczyszczenia treningu. W wariancie
z jądrem RBF trzeba dobrać $\nu$ oraz parametr $\gamma$, a także zadbać
o skalowanie kanałów o różnych jednostkach. Wyniki należy porównać
z Isolation Forest na identycznych przedziałach czasu, mierząc również
koszt inferencji na serwerze. HMI powinien prezentować wynik jako ocenę
zgodności z wyuczonym obszarem, bez nadawania mu znaczenia skalibrowanego
prawdopodobieństwa awarii urządzenia.

\subsubsection{Schmidl, Wenig i Papenbrock: porównanie detekcji czasowej}
Praca \emph{Anomaly Detection in Time Series: A Comprehensive Evaluation}
porównuje 71 algorytmów na 976 zbiorach szeregów czasowych. Autorzy
porządkują rodziny metod i analizują skuteczność, czas działania oraz
zależność wyników od właściwości danych. Rozróżniają m.in. anomalie punktowe,
nietypowe fragmenty przebiegów i anomalie związane z zależnościami między
kanałami. Udostępnione materiały uzupełniające pozwalają prześledzić
organizację porównania i wykorzystane implementacje~\cite{schmidl2022}.

Artykuł wybrano, ponieważ przeciwdziała ocenianiu algorytmu wyłącznie
na podstawie jednego udanego przykładu. Wynik konkretnej implementacji
zależy także od parametrów, reprezentacji danych i rodzaju zakłócenia.
Rozległy benchmark stanowi wzorzec rzetelnego porównania, ale jego ranking
nie rozstrzyga automatycznie, który model sprawdzi się w lokalnej sieci
ESP32. Projekt ma mniejszą skalę i własny zestaw kanałów pomiarowych.

Konkretny wkład to plan oceny Isolation Forest i One-Class SVM w kilku
scenariuszach, np. krótkiego wstrząsu, zmiany oświetlenia i utrzymującego
się odchylenia. Są to proponowane scenariusze projektu, a nie deklaracja
wykonania benchmarku autorów. Przy dostępnych etykietach należy raportować
precyzję, czułość i F1 oraz jasno określić, czy ocena dotyczy próbek,
czy całych zdarzeń. Bez etykiet nie należy przedstawiać udziału alarmów
jako skuteczności. Dla użytkownika HMI istotne będą dodatkowo liczba
fałszywych powiadomień i opóźnienie wykrycia, zaś dla serwera koszt
przetwarzania nowych pomiarów.

\subsection{Dokumentacje i rozwiązania techniczne}

Dokumentacje uzupełniają podstawy naukowe o informacje potrzebne do budowy
systemu. Obejmują implementację modeli, akwizycję i przesyłanie pomiarów
oraz rozwiązania stanowiące punkt odniesienia dla interfejsu HMI.

\subsubsection{scikit-learn: Novelty and Outlier Detection}
Dokumentacja rozróżnia wykrywanie obserwacji odstających w potencjalnie
zanieczyszczonym zbiorze oraz wykrywanie nowych nietypowych obserwacji
po uczeniu na danych odniesienia. Omawia m.in. Isolation Forest i One-Class
SVM oraz wspólny interfejs estymatorów. W predykcji etykieta $+1$ oznacza
obserwację uznaną za normalną, a $-1$ za odstającą; funkcja decyzyjna
pozwala pracować również z wynikiem ciągłym~\cite{sklearn_docs}.

Źródło wybrano jako bezpośrednią podstawę implementacji i interpretacji
wyników. Artykuły naukowe definiują metody, natomiast dokumentacja wyjaśnia
zachowanie udostępnionego oprogramowania. Jest to istotne zwłaszcza dla
Isolation Forest: parametr \texttt{contamination} wpływa na próg decyzji,
a nie dostarcza etykiet rzeczywistych zdarzeń. Dokumentacja wskazuje również
wrażliwość One-Class SVM na obserwacje odstające i potrzebę dostosowania
parametrów do zadania.

W projekcie źródło pozwala ujednolicić warstwę komunikacji modeli
z serwerem i HMI. Należy zapisywać zarówno wynik ciągły, jak i decyzję,
użyty próg oraz identyfikator modelu. Ujemny wynik funkcji decyzyjnej
powinien być interpretowany zgodnie z API, aby nie odwrócić znaczenia
alarmów. Wyników obu modeli nie należy traktować jako porównywalnych
prawdopodobieństw. W raporcie eksperymentalnym trzeba też podać rzeczywiście
użytą wersję biblioteki; odwołanie do dokumentacji \texttt{stable} samo
w sobie nie ustala wersji środowiska. Pozwoli to odtworzyć predykcje
i sprawdzić zgodność wyników prezentowanych na wykresach.

\subsubsection{Eclipse Mosquitto: broker MQTT}
Mosquitto jest otwartym brokerem implementującym MQTT, czyli komunikację
w modelu publikowania i subskrypcji. Nadawcy publikują wiadomości na
tematach, a broker przekazuje je zainteresowanym odbiorcom. Projekt
udostępnia również bibliotekę klienta i narzędzia \texttt{mosquitto\_pub}
oraz \texttt{mosquitto\_sub}. Dokumentacja i repozytorium opisują komponent
komunikacyjny, a nie gotowy system analizy pomiarów lub
detekcji anomalii~\cite{mosquitto_docs}.

Źródło wybrano, ponieważ broker stanowi naturalny element architektury
z kilkoma ESP32 i centralnym odbiorcą. Rozdzielenie publikujących urządzeń
od subskrybującego serwera ułatwia dodawanie kanałów oraz niezależne
testowanie odbioru telemetrii. Narzędzia konsolowe umożliwiają sprawdzenie
tematów i treści wiadomości przed uruchomieniem analizy ML. Użyteczność
Mosquitto nie oznacza, że sam broker zapewnia poprawność fizyczną danych
lub synchronizację pomiarów z różnych urządzeń.

Wkład w projekt to uzasadnienie warstwy transportowej i sposobu testowania
całego przepływu. Odbiorca powinien zachować temat MQTT, identyfikator
urządzenia, numer sekwencyjny i czas odbioru, aby umożliwić analizę przerw
lub powtórzeń. Dobór parametrów dostarczania i ponownego łączenia należy
udokumentować jako konfigurację projektu, bez zakładania, że przybycie
wiadomości oznacza jednoczesne próbkowanie. W HMI warto oddzielić brak
aktualnych danych od alarmu modelu: utrata połączenia ESP32 nie musi
generować nietypowego wektora, lecz może uniemożliwić jego utworzenie.
Broker i detektor rozwiązują zatem różne, uzupełniające się zadania.

\subsubsection{ThingsBoard: platforma IoT jako punkt odniesienia}
Dokumentacja ThingsBoard opisuje platformę obejmującą zarządzanie
urządzeniami, zbieranie telemetrii, pulpity wizualizacyjne oraz obsługę
alarmów i reguł przetwarzania. Materiały użytkownika pokazują, jak
powiązać dane z urządzeniami i prezentować ich stan w aplikacji.
Jest to przykład szerszego rozwiązania IoT, którego funkcje wykraczają
poza pojedynczy model detekcji~\cite{thingsboard_docs}.

Platformę wybrano jako punkt odniesienia dla zakresu funkcjonalnego HMI.
Porównanie z istniejącym produktem pomaga określić, co powinno być
dostępne dla operatora poza samym wykresem. Dokumentacja funkcji nie
stanowi jednak niezależnego pomiaru jakości platformy ani dowodu
skuteczności Isolation Forest lub One-Class SVM. Przy porównywaniu
rozwiązań trzeba wskazać konkretną edycję i konfigurację, ponieważ nie
wszystkie możliwości całego ekosystemu należą do jednego wariantu produktu.

Konkretnym wkładem w projekt jest zestaw kryteriów projektowania interfejsu:
identyfikacja węzła, aktualna wartość z jednostką, historia pomiarów,
czas wystąpienia alarmu i możliwość obejrzenia jego kontekstu.
Proponowane HMI powinno ponadto prezentować model i próg, które doprowadziły
do zgłoszenia, aby wspierać porównanie detektorów. Własny system można
ocenić względem ThingsBoard pod kątem tych funkcji, bez twierdzenia,
że przewyższa platformę komercyjną. Uzasadnieniem jego budowy jest
kontrolowany eksperyment z własnym przygotowaniem danych i algorytmami,
a nie konieczność odtworzenia wszystkich funkcji zarządzania
rozbudowaną infrastrukturą IoT.

\subsubsection{Grafana: wizualizacja i organizacja alarmów}
Dokumentacja Grafana Alerting opisuje reguły, które okresowo odpytują
źródła danych i sprawdzają warunki. Rozróżnia regułę alarmową, konkretne
instancje alarmu oraz kierowanie powiadomień do punktów kontaktowych.
Powiadomienia można organizować bezpośrednio lub za pomocą polityk.
Pozwala to oddzielić obliczanie warunku od sposobu informowania
użytkownika~\cite{grafana_docs}.

Źródło wybrano jako praktyczny wzorzec warstwy alarmowania i punkt
odniesienia dla HMI. Samo wykrycie odchylenia przez model nie określa
jeszcze, kiedy wysłać wiadomość i jak uniknąć wielu zgłoszeń dotyczących
tego samego zdarzenia. Dokumentacja wyjaśnia te pojęcia na poziomie
systemowym, ale nie jest publikacją potwierdzającą skuteczność
algorytmów ML. Przedstawione mechanizmy nie oznaczają również, że
Grafana jest już elementem implementacji projektu.

W projekcie proponuje się wykorzystać ten podział do zaprojektowania
osobnego rejestru wyników i rejestru zdarzeń. Kolejne alarmowe próbki
z jednego okresu zakłócenia mogą być prezentowane operatorowi jako
jedno zdarzenie z czasem rozpoczęcia i zakończenia, zamiast wielu
niezależnych wiadomości. Jest to propozycja dla własnego systemu,
wymagająca określenia reguł łączenia zdarzeń. Na wspólnej osi czasu
HMI powinien pokazywać pomiary, wynik detektora i próg. Przerwa
w telemetrii wymaga osobnej obsługi, ponieważ brak danych nie
jest równoznaczny z normalną pracą. Taka organizacja ułatwi ocenę
użyteczności alarmów oraz powiązanie ich z eksperymentami fizycznymi.

\subsubsection{Espressif: biblioteki Arduino-ESP32}
Oficjalna dokumentacja Arduino-ESP32 opisuje biblioteki udostępniające
funkcje układów ESP32 przez interfejsy Arduino. Zestawienie obejmuje
m.in. GPIO, I2C, SPI, timery i komunikację Wi-Fi oraz pokazuje różnice
w dostępności funkcji między układami rodziny. Materiały producenta
są podstawą sprawdzania możliwości wybranego mikrokontrolera
i korzystania z jego interfejsów~\cite{espressif_docs}.

Źródło wybrano ze względu na zgodność z warstwą sprzętową projektu.
Trzy węzły muszą regularnie odczytywać czujniki i przekazywać
pomiary do serwera, dlatego sposób obsługi magistral i połączenia
sieciowego wpływa na jakość danych wejściowych ML. Dokumentacja
rdzenia Arduino nie zastępuje jednak dokumentacji konkretnych
czujników AHT20, BMP280, BH1750 i ADXL345 ani dokumentacji
wybranej biblioteki klienta MQTT. Nie należy przypisywać jej
gotowych sterowników lub algorytmów, których nie opisuje.

Konkretnym wkładem jest możliwość sprawdzenia interfejsów i ograniczeń
sprzętowych przed interpretacją nietypowych odczytów. W oprogramowaniu
węzłów należy jednoznacznie określić jednostki, częstotliwość raportowania
i znaczenie znaczników czasu, a dla akcelerometru również sposób
wyznaczania RMS i maksimum. Są to wymagania projektowe wynikające
z potrzeby porównywalności pomiarów. Numer sekwencyjny i czas urządzenia
pomogą rozpoznać restart lub lukę w transmisji; czas odbioru na serwerze
pozostaje osobną informacją. Tak przygotowana telemetria ułatwi
oddzielenie problemów akwizycji od alarmów Isolation Forest i One-Class SVM
oraz prezentację wiarygodnego stanu węzłów w HMI.

\subsection{Wnioski dla projektu}
Literatura uzasadnia porównanie dwóch mechanizmów detekcji na wspólnej
reprezentacji danych. W metodologii należy rozdzielić przygotowanie pomiarów,
uczenie, wybór progu i obsługę alarmów. Skuteczność trzeba ocenić na danych
projektu; nie wynika ona automatycznie z cytowania publikacji.

\section{Metodologia zbierania i podziału danych}
Pomiary z trzech węzłów ESP32 zapisano na serwerze i zestawiono ze sobą na podstawie czasu odbioru. Dane podzielono chronologicznie na trzy zbiory w proporcji 70/15/15: najwcześniejsze pomiary przeznaczono na trening, kolejne na walidację, a najnowsze na test.

\subsection{Zbieranie danych}
W systemie działają trzy węzły ESP32: pierwszy z czujnikami AHT20 i BMP280, drugi z czujnikiem oświetlenia BH1750, a trzeci z akcelerometrem ADXL345. Telemetria trafia przez MQTT na serwer i jest zapisywana w pliku CSV. W repozytorium plik \texttt{data/raw/telemetry.csv} zawiera czas odbioru \texttt{received\_at}, temat MQTT, identyfikator urządzenia, nazwę kanału, czas urządzenia \texttt{ts\_ms}, numer sekwencyjny, wartość oraz jednostkę. Każdy wiersz opisuje jeden kanał, a nie pełny wektor wszystkich pomiarów.

Zebrany zbiór obejmuje około 39 godzin pomiarów.

\subsection{Synchronizacja i reprezentacja pomiarów}
Dane z poszczególnych kanałów sprowadzono do wspólnej siatki czasowej o kroku 5 s. Jako punkt odniesienia przyjęto czas odbioru na serwerze, dzięki czemu możliwe było zestawienie kanałów z różnych urządzeń. Nie jest to jednak równoznaczne z jednoczesnym próbkowaniem: opóźnienia w transmisji i różne częstotliwości raportowania nadal wprowadzają pewną niepewność.

Każda przygotowana próbka składa się z sześciu cech:
\begin{enumerate}
	\item \texttt{AHT20\_temperature} -- temperatura;
	\item \texttt{AHT20\_humidity} -- wilgotność względna;
	\item \texttt{BMP280\_pressure} -- ciśnienie atmosferyczne;
	\item \texttt{BH1750} -- natężenie oświetlenia;
	\item \texttt{ADXL345\_accel\_rms} -- raportowana miara RMS przyspieszenia;
	\item \texttt{ADXL345\_accel\_max} -- raportowana maksymalna wartość przyspieszenia.
\end{enumerate}
Znacznik \texttt{received\_at} służy wyłącznie do uporządkowania danych i nie jest cechą modelu. Obie miary z akcelerometru potraktowano jako osobne kanały.

\subsection{Podział chronologiczny i oczyszczanie}
Obserwacje podzielono chronologicznie w proporcji 70/15/15: najwcześniejsza część jest przeznaczona do uczenia, następna do walidacji, a ostatnia do końcowego testu. Nie zastosowano losowego mieszania próbek. Sąsiednie odczyty są ze sobą powiązane w czasie, więc przy losowym podziale bardzo podobne fragmenty przebiegu mogłyby trafić jednocześnie do zbioru treningowego i do oceny. Podział chronologiczny lepiej odpowiada sytuacji, w której model analizuje pomiary z późniejszego okresu, choć nie eliminuje całkowicie zależności między zbiorami.

Zbiór walidacyjny będzie służył do obserwowania zachowania modelu i doboru parametrów. Zbiór testowy pozostaje zarezerwowany do oceny końcowej, po podjęciu tych decyzji, i nie powinien służyć do dobierania progu. Podane proporcje dotyczą samego podziału chronologicznego; po przygotowaniu i oczyszczeniu danych liczebności zbiorów mogą się od nich nieznacznie różnić.

Potencjalne wstrząsy usuwa się wyłącznie ze zbioru treningowego, aby model nie uczył się zakłóceń jako zachowania referencyjnego. Ze zbiorów walidacyjnego i testowego nie usuwa się obserwacji tylko dlatego, że wyglądają nietypowo, ponieważ sztucznie ułatwiłoby to ocenę. Regułę oczyszczania należy ustalić na podstawie zbioru treningowego i wiedzy o eksperymencie, bez korzystania z późniejszych wyników. Nie stanowi ona jednak potwierdzenia etykiet anomalii.

Liczebności zbiorów po przygotowaniu danych zestawiono w tabeli. Dotyczą one gotowych wektorów sześciu cech, a nie surowych wierszy telemetrii.
\begin{center}
	\begin{tabular}{lrr}
		\hline
		Zbiór & Liczba próbek & Liczba cech \\
		\hline
		Treningowy & 19\,634 & 6 \\
		Walidacyjny & 4\,222 & 6 \\
		Testowy & 4\,222 & 6 \\
		\hline
	\end{tabular}
\end{center}

\clearpage
\begingroup
\small
\bibliographystyle{unsrt}
\bibliography{bibliografia}
\endgroup

\end{document}
